\documentclass[12pt,a4paper]{article}
\usepackage[T1]{fontenc}
\usepackage[utf8]{inputenc}
\usepackage[ngerman]{babel}
\usepackage{amsmath,amssymb,amsthm}
\usepackage{geometry}
\geometry{margin=2.5cm}
\usepackage{fancyhdr}
\pagestyle{fancy}
\fancyhf{}
\fancyhead[C]{\small Lehrbuch der Algebra --- Band II, Zwölfter Abschnitt}
\fancyfoot[C]{\thepage}
\renewcommand{\headrulewidth}{0.4pt}

\theoremstyle{definition}
\newtheorem*{satz}{Satz}

\title{Lehrbuch der Algebra --- Band 1, Teil 8}
\author{Heinrich Weber}
\date{}

\begin{document}
\maketitle
\thispagestyle{empty}

\begin{center}
\Large\textbf{Zwölfter Abschnitt.}

\medskip
\Large\textbf{Die Doppeltangenten einer Curve vierter Ordnung.}
\end{center}

\bigskip

\section*{\S. 95.}
\addcontentsline{toc}{section}{\S. 95. Anzahl der Doppeltangenten einer Curve vierter Ordnung.}

\begin{center}
\textbf{Anzahl der Doppeltangenten einer Curve vierter Ordnung.}
\end{center}

Ein geometrisches Problem, das wegen seiner mannigfachen
Beziehungen zu anderen Gebieten von besonderer Bedeutung ist,
was zugleich in ähnlicher Weise, wie das Problem der Wende-
punkte der Curven dritter Ordnung, merkwürdige algebraische
Verhältnisse bietet, ist das der Doppeltangenten der Curven
vierter Ordnung.

Unter einer \textit{Doppeltangente} einer Curve versteht man,
wie schon im \S.~86 bemerkt ist, eine gerade Linie, die die Curve
in zwei verschiedenen Punkten berührt. Bei Curven von niedri-
gerer als der vierten Ordnung können Doppeltangenten nicht
auftreten.

Bei den Curven vierter Ordnung hat eine Doppeltangente
ausser den Berührungspunkten keinen Punkt mit der Curve
gemein. In besonderen Fällen können die beiden Berührungs-
punkte auch zusammenfallen. Dann haben wir Linien mit vier-
punktiger Berührung.

Die Bestimmung der Doppeltangenten wird als algebraisches
Problem von einer gewissen algebraischen Gleichung abhängen,
deren Grad gleich der Anzahl der Doppeltangenten ist, und die
erste Frage ist die nach dem Grade dieser Gleichung, also nach
der \textit{Anzahl} der Doppeltangenten. Wir beschränken uns hier auf
die Betrachtung von Curven vierter Ordnung, obwohl der Weg,
den wir gehen, auch auf Curven höherer Ordnung anwendbar
ist. Wir setzen auch voraus, dass die Curve vierter Ordnung
frei von singulären Punkten sei\footnote{Jacobi, ``Ueber die Anzahl der Doppeltangenten ebener algebraischer Curven'', Crelle's Journal, Bd.~40 (1850). Clebsch, ``Bemerkung zu Jacobi's Beweis für die Anzahl der Doppeltangenten'', ebend., Bd.~68 (1864). Aus der ziemlich umfassenden Literatur über die Theorie der Doppeltangenten einer Curve vierter Ordnung sind die folgenden Schriften hervorzuheben. Zunächst mehrere Arbeiten von Hesse in Crelle's Journal, Bd.~41, 49, 52. Ferner Steiner, Eigenschaften der Curven vierter Ordnung rücksichtlich ihrer Doppeltangenten, Crelle's Journal, Bd.~49. Aronhold, Monatsber. d. Berl. Akademie 1864. Geiser, ``Ueber die Doppeltangenten einer ebenen Curve $4^{\text{ten}}$ Grades'', Math. Annalen, Bd.~1 (1869). Cayley, Crelle's Journal, Bd.~68 (1868). Auch das oben citirte Werk von Salmon ist hier hervorzuheben. In neuerer Zeit wurde die Theorie der Doppeltangenten im Zusammenhange mit der Theorie der Abel'schen Functionen weiter ausgebildet. Riemann, Gesammelte mathematische Werke (2.~Aufl., 1892) XXXI, aus dem Nachlass. Weber, Theorie der Abel'schen Functionen vom Geschlecht~3, Berlin 1876. Frobenius, Crelle's Journal, Bd.~99 u.~103. Die algebraische Seite der Frage ist in zwei Abhandlungen: Nöther, Math. Annalen, Bd.~15 und Weber, ebend., Bd.~23 behandelt.}.

Es möge jetzt $f(x_1, x_2, x_3)$ oder kürzer $f(x)$ eine ternäre
Form $4^{\text{ten}}$ Grades sein, die, gleich Null gesetzt, eine Curve vierter
Ordnung ohne singulären Punkt darstellt, die wir die Curve~$f$
nennen.

Setzen wir in dieser Gleichung
\begin{equation}\label{eq:1}
f(x + ty) = 0
\end{equation}
an Stelle der Variablen $x_1, x_2, x_3$ Ausdrücke von der Form
\[x_1 + ty_1,\ x_2 + ty_2,\ x_3 + ty_3,\]
so ergiebt sich eine Gleichung $4^{\text{ten}}$ Grades in Bezug auf $t$:
\begin{equation}\label{eq:2}
f(x + ty) = 0,
\end{equation}
deren Wurzeln, wenn $(x)$ und $(y)$ zwei feste Punkte sind, in
$x + ty$ eingesetzt, die Coordinaten der Schnittpunkte der Ver-
bindungslinie der Punkte $(x)$, $(y)$, die wir als die Linie $(x,y)$
bezeichnen wollen, mit der Curve~$f$ geben.

Es werde nun die Function $f(x + ty)$ nach Potenzen von~$t$
geordnet. Dies giebt (nach Bd.~I, \S.~60):
\begin{multline}\label{eq:3}
f(x + ty) = P_0(x,y) + 4tP_1(x,y) + 6t^2 P_2(x,y)\\
+ 4t^3 P_3(x,y) + t^4 P_4(x,y),
\end{multline}
worin die $P_\nu(x,y)$ die Polaren von $f(x)$ sind, also homogene
Functionen $\nu^{\text{ter}}$ Ordnung in Bezug auf $y$, und $(4-\nu)^{\text{ter}}$ Ordnung
in Bezug auf $x$. Es ist nämlich
\begin{equation}\label{eq:4}
\begin{aligned}
P_0(x,y) &= f(x),\\
P_1(x,y) &= \tfrac{1}{4} \sum y_i f'(x_i),\\
P_2(x,y) &= \tfrac{1}{6} \sum y_i y_k f''(x_i x_k),\\
P_3(x,y) &= \tfrac{1}{4} \sum y_i f'(y_i),\\
P_4(x,y) &= f(y).
\end{aligned}
\end{equation}

Wenn $x$ auf der Curve $f$ liegt, so wird $P_0 = 0$, und eine
Wurzel der Gleichung~(\ref{eq:2}) verschwindet. Nehmen wir ausser-
dem noch $(y)$ auf der Tangente im Punkte $(x)$ an, so ist auch
$P_1(x,y) = 0$, und es verschwinden zwei Wurzeln von~(\ref{eq:2}). Die
beiden übrigen werden nach~(\ref{eq:3}) durch die quadratische Gleichung
\begin{equation}\label{eq:5}
6P_2 + 4tP_3 + t^2 P_4 = 0
\end{equation}
bestimmt. Wenn diese Gleichung zwei gleiche Wurzeln hat, so
wird die Linie $(x,y)$ noch einen zweiten Berührungspunkt mit
der Curve~$f$ haben, d.\,h. sie wird Doppeltangente sein. Die Be-
dingung hierfür ist aber die, dass die Discriminante der Glei-
chung~(\ref{eq:5}) verschwindet, oder dass
\begin{equation}\label{eq:6}
R(x,y) = 2P_3^2 - 3P_2 P_4
\end{equation}
gleich Null sei. Hierin ist $R(x,y)$ eine in Bezug auf jedes der
beiden Systeme $(x)$ und $(y)$ homogene Function, und zwar für
die $x$ vom $2^{\text{ten}}$, für die $y$ vom $6^{\text{ten}}$ Grade. Es ist aber noch zu
bemerken, dass die Gleichung $R = 0$ auch dann erfüllt ist,
wenn $(x)$ ein beliebiger Punkt der Curve ist, und $(y)$ mit $(x)$
zusammenfällt, weil dann
\[P_2(x,x) = P_3(x,x) = P_4(x,x) = f(x) = 0\]
wird. Sonst aber wird $R$ nur dann verschwinden, wenn die Linie
$(x,y)$ eine Doppeltangente ist. Wir stellen also den Satz auf:

\begin{satz}
Die Gleichung $R(x,y) = 0$ ist, wenn $(x)$ ein Punkt
der Curve~$f$ und $(y)$ ein von $(x)$ verschiedener
Punkt der Tangente in $(x)$ ist, die nothwendige
und hinreichende Bedingung dafür, dass $(x)$ ein
Berührungspunkt einer Doppeltangente sei.
\end{satz}

Es kommt nun darauf an, aus dieser Bedingung eine andere
herzuleiten, die nur die $x$ allein enthält, und die für keine
anderen Punkte als die Berührungspunkte der Doppeltangenten
befriedigt ist.

Die Variablen $y$ genügen der Tangentengleichung
\begin{equation}\label{eq:7}
y_1 f'(x_1) + y_2 f'(x_2) + y_3 f'(x_3) = 0.
\end{equation}

Bezeichnen wir mit $b_1, b_2, b_3$ irgend drei willkürliche Con-
stanten, und setzen
\begin{equation}\label{eq:8}
b_x = b_1 x_1 + b_2 x_2 + b_3 x_3,
\end{equation}
so dass $b_x = 0$ die Gleichung einer willkürlichen geraden Linie
ist, so können wir zu~(\ref{eq:7}) noch die Gleichung
\begin{equation}\label{eq:9}
b_1 y_1 + b_2 y_2 + b_3 y_3 = 0
\end{equation}
hinzunehmen, also $(y)$ als den Schnittpunkt der Tangente in $(x)$
mit der beliebigen geraden Linie $b_x$ auffassen. Dann erhalten wir
\begin{equation}\label{eq:10}
\begin{aligned}
\rho y_1 &= b_3 f'(x_2) - b_2 f'(x_3),\\
\rho y_2 &= b_1 f'(x_3) - b_3 f'(x_1),\\
\rho y_3 &= b_2 f'(x_1) - b_1 f'(x_2),
\end{aligned}
\end{equation}
und hierdurch ist die Bedingung~(\ref{eq:7}) identisch befriedigt. Durch
diese Substitution geht $R(x,y)$ in eine homogene Function
$20^{\text{ten}}$ Grades der $x$ und $6^{\text{ten}}$ Grades der $b$ über, die wir mit
$D(x,b)$ bezeichnen wollen. Diese Function $D(x,b)$ verschwindet
aber ausser in den Berührungspunkten der Doppeltangenten noch
in den vier Schnittpunkten der Geraden $b_x$ mit der Curve~$f$.
Diese Bedingung muss nun noch so umgeformt werden, dass sie
von den $b$ unabhängig wird.

Gehen wir von dem Punkte $(y)$ zu einem beliebigen anderen
Punkte der Tangente über, so können wir dies dadurch er-
reichen, dass wir $y$ durch $x + \lambda y$ ersetzen, worin $\lambda$, $\mu$ zwei
willkürliche Parameter bedeuten. Dadurch geht $f(x + ty)$ in
\[f\bigl(x + t(x + \lambda y)\bigr) = f\bigl((1+t)x + t\lambda y\bigr)\]
über, und die Entwickelung~(\ref{eq:3}) ergiebt, wenn wir
\[P_\nu(x,y),\ P_{\nu+1}(x,y),\ P_{\nu+2}(x,\lambda x + \mu y)\]
gleich Null setzen:
\begin{multline*}
6t^2(1+\lambda t)^2 \mu^2 P_2(x,y) + 4t^3(1+\lambda t)\mu P_3(x,y) + t^4 P_4(x,y)\\
= 6P_2(x,\mu x + \lambda y) + 4tP_3(x,\mu x + \lambda y) + t^2 P_4(x,\mu x + \lambda y).
\end{multline*}

\newpage
\section*{\S. 96.}
\addcontentsline{toc}{section}{\S. 96. Die Hesse'sche Curve.}

\begin{center}
\textbf{Die Hesse'sche Curve.}
\end{center}

\end{document}
